% ======================================================================
% Capítulo: Diseño del sistema
% Fase 2 del anteproyecto. Extensión orientativa: 8--10 páginas.
% ======================================================================
\chapter{Diseño del sistema}
\label{cap:diseno}

\section{Arquitectura general}
\label{sec:arquitectura}
% TODO: Los tres componentes desacoplados (Go, .NET, Angular), gateway, workers y almacenes.

\section{Arquitectura limpia en .NET}
\label{sec:clean}
% TODO: Capas, puertos y adaptadores, tests de arquitectura.

\section{Diagramas UML}
\label{sec:uml}
% TODO: Clases de los módulos clave y secuencias de pregunta en el chat e ingesta de un boletín.

\section{Modelo de datos híbrido}
\label{sec:modelo-datos}
% TODO: Diagrama E-R de PostgreSQL y colecciones de Qdrant (\texttt{document\_chunks}, \texttt{child\_fragments}).

\section{Indexación vectorial}
\label{sec:indexacion}
% TODO: Modelos de embeddings y dimensiones, similitud coseno, HNSW, límite de pgvector y migración a Qdrant (ADR 0001).

\section{Flujo de ingesta asíncrona}
\label{sec:ingesta-diseno}
% TODO: Contrato de RabbitMQ, colas de reintento y DLQ, idempotencia y estados del documento.

\section{Pipelines RAG}
\label{sec:pipelines}
% TODO: Baseline (chunks fijos, búsqueda densa) frente a jerárquica (expansión de consulta, híbrida RRF, contexto padre, lenguaje claro).

\section{Abstracción del proveedor de IA y control de coste}
\label{sec:proveedor-ia}
% TODO: Semantic Kernel, fábrica de modelos por tenant, fallback ante 429 y registro de coste por llamada.

\section{Tiempo real y seguridad}
\label{sec:seguridad}
% TODO: SignalR, multi-tenancy con Auth0 (acceso por ayuntamiento), roles y cifrado de claves.
